\documentclass[11pt]{article}
\usepackage[letterpaper,margin=0.72in,top=0.82in,bottom=0.72in,headheight=22pt]{geometry}
\usepackage[T1]{fontenc}
\usepackage{lmodern}
\usepackage{microtype}
\usepackage[table]{xcolor}
\usepackage{amsmath,amssymb,mathtools}
\usepackage{array,longtable,booktabs,tabularx}
\usepackage[inline]{enumitem}
\usepackage[most]{tcolorbox}
\usepackage{fancyhdr}
\usepackage{hyperref}

\definecolor{Primary}{HTML}{17324D}
\definecolor{Accent}{HTML}{2C7A7B}
\definecolor{CueTint}{HTML}{EAF1F5}
\definecolor{NoteTint}{HTML}{F8FAFB}
\definecolor{Rule}{HTML}{B8C6CF}
\hypersetup{hidelinks,pdftitle={Chapter 17: Strings and Evolutionary Trees},pdfsubject={Cornell notes}}
\setlist[itemize]{leftmargin=1.25em,itemsep=1pt,topsep=1.5pt,parsep=0pt}
\setlength{\parindent}{0pt}
\setlength{\parskip}{3pt}
\renewcommand{\arraystretch}{1.16}
\emergencystretch=2em

\pagestyle{fancy}
\fancyhf{}
\fancyhead[L]{\sffamily\small\color{Primary} Chapter 17}
\fancyhead[R]{\sffamily\small\color{Primary} Evolutionary trees}
\fancyfoot[C]{\sffamily\small\color{Primary}\thepage}
\renewcommand{\headrulewidth}{0.5pt}
\renewcommand{\headrule}{\hbox to\headwidth{\color{Accent}\leaders\hrule height \headrulewidth\hfill}}

\newcolumntype{C}{>{\raggedright\arraybackslash}p{0.275\textwidth}}
\newcolumntype{N}{>{\raggedright\arraybackslash}p{0.655\textwidth}}
\newcommand{\cnrow}[2]{%
  \cellcolor{CueTint}\textbf{\sffamily\color{Primary}#1} &
  \cellcolor{NoteTint}#2 \\[4pt]\arrayrulecolor{Rule}\hline}

\begin{document}

\begin{tcolorbox}[
  enhanced,breakable,colback=Primary!4,colframe=Primary,boxrule=0.9pt,
  arc=2mm,left=3mm,right=3mm,top=2.5mm,bottom=2.5mm]
{\sffamily\color{Accent}\bfseries CHAPTER 17}\\[-1pt]
{\sffamily\LARGE\bfseries\color{Primary} Strings and Evolutionary Trees}

\vspace{2pt}
\textbf{Purpose.} Connect pairwise sequence distances and character changes to ultrametric, additive, parsimony, and perfect-phylogeny tree reconstruction.

\textbf{Learning objectives}
\begin{itemize}
  \item Recognize ultrametric and additive distance conditions.
  \item Reconstruct or validate a weighted tree from exact distances.
  \item Compute a parsimony score on a fixed tree.
  \item Explain the mutual dependence of alignment and phylogeny.
\end{itemize}

\textbf{Prerequisites.} Weighted trees, metrics, dynamic programming, sequence alignment, and basic evolutionary concepts.
\end{tcolorbox}

\begin{longtable}{@{}C!{\color{Accent}\vrule width 0.8pt}N@{}}
\rowcolor{Primary}
\color{white}\sffamily\bfseries Cue / question &
\color{white}\sffamily\bfseries Notes \\ \hline
\endfirsthead
\rowcolor{Primary}
\color{white}\sffamily\bfseries Cue / question &
\color{white}\sffamily\bfseries Notes (continued) \\ \hline
\endhead
\cnrow{What are distance- and character-based methods?}{Distance methods compress each sequence pair into one dissimilarity value and fit a tree to the matrix. Character methods retain site states and score changes along tree edges. Compression improves efficiency but can discard which sites support each branch.}
\cnrow{\S17.1: What is an ultrametric?}{An ultrametric satisfies $d(x,z)\le\max\{d(x,y),d(y,z)\}$ for every triple. Equivalently, the two largest distances in a triple are equal. It corresponds to a rooted tree whose leaves are equally distant in time from the root under a strict clock model.}
\cnrow{How is an ultrametric tree reconstructed?}{Repeatedly cluster the closest compatible taxa or distance classes and assign an ancestor height consistent with their distance. After joining a cluster, update its distance to other clusters according to the ultrametric rule. Exact input must satisfy the triple condition; noisy data require an approximation method.}
\cnrow{What is the correctness cue for clustering?}{In an exact ultrametric, a closest pair forms a valid lowest cluster because no third leaf can be closer to one member without violating the triple condition. Contracting that cluster preserves ultrametricity on the reduced instance.}
\cnrow{\S17.2: What is an additive distance?}{A distance matrix is additive if there is a weighted unrooted tree whose path lengths equal all pairwise distances. The four-point condition states that, for any four leaves, the two largest of the three pair-sums are equal. Unlike ultrametrics, leaves need not be equidistant from a root.}
\cnrow{How is an additive tree recovered?}{Identify cherries or infer limb lengths from distance differences, remove or contract a leaf, recursively reconstruct the smaller tree, and attach the leaf at the uniquely implied point. Nonnegative edge lengths and consistency checks are essential.}
\cnrow{\S17.3: What is parsimony?}{Assign states to internal nodes so the total number or cost of changes over edges is minimized. On a fixed tree, a bottom-up DP stores, for every node and possible character, the best cost for its subtree; a top-down pass recovers ancestral assignments.}
\cnrow{What is the fixed-tree parsimony recurrence?}{For node $v$ and state $a$, sum over children $u$ the minimum of $F_u(b)+c(a,b)$ over child states $b$. Leaves permit only their observed state. Columns can be scored independently when the model has no interaction across sites.}
\cnrow{Why is choosing the tree harder?}{The fixed-tree DP is polynomial, but searching tree topologies grows superexponentially and common optimization variants are computationally hard. Branch-and-bound, heuristics, and consensus across good trees are therefore common.}
\cnrow{\S17.4: Why is the ultrametric problem central?}{Ultrametric reconstruction offers a clean example in which local distance conditions completely characterize a rooted tree. It clarifies how assumptions such as a molecular clock create algorithmic tractability and how violations signal model mismatch or noise.}
\cnrow{\S17.5: How are parsimony and Steiner trees related?}{Treat possible ancestral states as Steiner vertices connecting observed leaves in a metric state space. Minimizing evolutionary change resembles a Steiner-tree problem. Perfect phylogeny is the special case where each character state arises without repeated independent changes under its model.}
\cnrow{What is perfect phylogeny testing?}{Encode each character's state partition and ask whether one tree realizes all characters without forbidden recurrence. Compatibility conditions vary with binary, multistate, ordered, or unrooted models; the character model must be stated before applying a test.}
\cnrow{\S17.6: What is phylogenetic alignment?}{Instead of treating the alignment as fixed, score insertions, deletions, and substitutions along a tree and seek ancestral sequences plus alignments on edges. This models evolution more directly but combines already difficult alignment and topology choices.}
\cnrow{\S17.7: Why do alignment and tree construction interact?}{An alignment defines the columns used to infer a tree, while a tree guides which residues should be aligned as homologous. Alternating improvement can stabilize at a local optimum, so robustness should be checked across guide trees, scoring choices, and data subsets.}
\cnrow{\S17.8: What should practice emphasize?}{Test triple and four-point conditions, reconstruct a small exact tree, run fixed-tree parsimony by hand, and identify which evolutionary assumptions are embedded in each objective.}
\end{longtable}


\begin{tcolorbox}[enhanced,breakable,title={Chapter synthesis},
  colback=Accent!5,colframe=Accent,fonttitle=\small\sffamily\bfseries,fontupper=\footnotesize,
  boxsep=0.6mm,left=1.5mm,right=1.5mm,top=1mm,bottom=1mm,before skip=3pt,after skip=3pt]
Evolutionary reconstruction can start from distances or from characters. Ultrametrics encode rooted clocklike trees through a triple condition; additive metrics encode weighted unrooted trees through a four-point condition; parsimony assigns internal character states by tree DP. When topology and alignment are also unknown, the problems become mutually dependent and computationally harder.
\end{tcolorbox}

\begin{tcolorbox}[enhanced,breakable,title={Key takeaways},
  colback=white,colframe=Primary,fonttitle=\small\sffamily\bfseries,fontupper=\footnotesize,
  boxsep=0.6mm,left=1.5mm,right=1.5mm,top=1mm,bottom=1mm,before skip=3pt,after skip=3pt]
\begin{itemize}
  \item An ultrametric's two largest distances in every triple are equal.
  \item An additive tree metric satisfies the four-point condition.
  \item Exact distance conditions permit recursive reconstruction.
  \item Parsimony on a fixed tree is a bottom-up state DP.
  \item Choosing a topology is much harder than scoring one.
  \item Perfect phylogeny depends on the character-state model.
  \item Alignment and phylogeny influence one another.
\end{itemize}
\end{tcolorbox}

\begin{tcolorbox}[enhanced,breakable,title={Algorithm selection / when to use this},
  colback=Primary!3,colframe=Primary,fonttitle=\small\sffamily\bfseries,fontupper=\footnotesize,
  boxsep=0.6mm,left=1.5mm,right=1.5mm,top=1mm,bottom=1mm,before skip=3pt,after skip=3pt]
\begin{itemize}
  \item Use ultrametric methods when a rooted clock model is credible.
  \item Use additive-distance reconstruction for exact or well-fitted path metrics without equal root-to-leaf times.
  \item Use fixed-tree parsimony to infer economical ancestral states.
  \item Use joint or iterative methods cautiously when alignment and topology are both uncertain.
\end{itemize}
\end{tcolorbox}

\begin{tcolorbox}[enhanced,breakable,title={Self-test questions},
  colback=white,colframe=Accent,fonttitle=\small\sffamily\bfseries,fontupper=\footnotesize,
  boxsep=0.6mm,left=1.5mm,right=1.5mm,top=1mm,bottom=1mm,before skip=3pt,after skip=3pt]
\begin{itemize}
  \item State and interpret the ultrametric triple condition.
  \item State the additive four-point condition.
  \item Why can the closest pair be contracted in an exact ultrametric?
  \item Write the fixed-tree parsimony recurrence.
  \item How is perfect phylogeny different from ordinary minimum parsimony?
  \item Why can changing an alignment change the inferred tree?
\end{itemize}
\end{tcolorbox}

{\footnotesize
\textbf{Terms and notation to remember.} ultrametric, molecular clock, additive metric, four-point condition, cherry, limb length, parsimony, ancestral state, Steiner tree, perfect phylogeny.

\textbf{Connections.} Pairwise and multiple alignments supply the distances or characters used here; genome rearrangement distances broaden evolutionary comparison beyond point edits.
}

\end{document}
